\documentclass[11pt]{article}
\usepackage[T1]{fontenc}
\usepackage[utf8]{inputenc}
\usepackage{lmodern}
\usepackage[a4paper,margin=1in]{geometry}
\usepackage{microtype}
\usepackage{amsmath,amssymb,amsthm,mathtools}
\usepackage{booktabs,array}
\usepackage[dvipsnames]{xcolor}
\usepackage[colorlinks=true,linkcolor=MidnightBlue,citecolor=MidnightBlue,
            urlcolor=MidnightBlue]{hyperref}
\usepackage{enumitem}
\usepackage{setspace}
\onehalfspacing
\setlist[itemize]{topsep=4pt,itemsep=2pt,leftmargin=1.5em}
\setlist[enumerate]{topsep=4pt,itemsep=2pt,leftmargin=1.8em}

\newtheorem{theorem}{Theorem}[section]
\newtheorem{lemma}[theorem]{Lemma}
\newtheorem{proposition}[theorem]{Proposition}
\newtheorem{corollary}[theorem]{Corollary}
\theoremstyle{definition}
\newtheorem{definition}[theorem]{Definition}
\newtheorem{construction}[theorem]{Construction}
\theoremstyle{remark}
\newtheorem{remark}[theorem]{Remark}

\newcommand{\DR}{\mathsf{DR}}
\newcommand{\PO}{\mathsf{PO}}
\newcommand{\PP}{\mathsf{PP}}
\newcommand{\PPI}{\mathsf{PPI}}
\newcommand{\EQ}{\mathsf{EQ}}
\newcommand{\RCC}{\{\DR,\PO,\PP,\PPI\}}
\newcommand{\cl}{\operatorname{cl}}
\newcommand{\Tp}{\operatorname{Tp}}
\newcommand{\tp}{\operatorname{tp}}
\newcommand{\md}{\operatorname{md}}
\newcommand{\st}{\operatorname{st}}
\newcommand{\comp}{\operatorname{comp}}
\newcommand{\inv}{\operatorname{inv}}
\newcommand{\orig}{\operatorname{orig}}
\newcommand{\Need}{\operatorname{Need}}
\newcommand{\Safe}{\operatorname{Safe}}
\newcommand{\Core}{\mathsf{Core}}
\newcommand{\Out}{\mathsf{Out}}
\newcommand{\Front}{\mathsf{Front}}
\newcommand{\Wit}{\operatorname{Wit}}
\newcommand{\Iface}{\operatorname{Iface}}
\newcommand{\Child}{\operatorname{Child}}
\newcommand{\Par}{\operatorname{Par}}
\newcommand{\Dom}{\mathsf{Dom}}
\newcommand{\ALCIRCC}[1]{\mathcal{ALCI}_{\mathsf{RCC#1}}}

\title{\textbf{Completeness Extraction for $\ALCIRCC{5}$}\\[0.4em]
\large Closing the model-to-quotient gap in the split-forest approach}
\author{Claude (Anthropic, Opus 4.6)\\
\small Prompted by Michael Wessel}
\date{April 2026}

\begin{document}
\maketitle

\begin{center}
\fbox{\parbox{0.92\textwidth}{%
\textbf{Disclaimer.}\;
This document was generated by Claude (Anthropic), prompted by Michael
Wessel.  The results have not been independently verified.  We invite
scrutiny and corrections.}}
\end{center}

\begin{abstract}
The split-forest approach to $\ALCIRCC{5}$ decidability~\cite{SplitForest}
proves the \emph{soundness} direction in full detail: every valid finite
rank-$d$ quotient unfolds into a genuine strong-$\EQ$ model via canonical
refinements of disjunctive RCC5 networks.  The \emph{completeness}
direction---that every satisfiable concept admits a valid finite
quotient---is stated as a proof sketch, with the extraction of
support-closed descriptor tables and relation-sensitive witness-menu
filters from a model left compressed.

We write out this extraction in full.  The key observation is that
\textbf{every relation realized in a valid model is already
relation-safe for the types of its endpoints}, because the model
satisfies all universal constraints.  When the model is quotiented by
rank-$d$ state equality, the disjunctive domains consist entirely of
relation-safe labels, so $\Need_R$-filtering preserves non-emptiness
trivially.  Combined with the composition-consistency of the model's
atomic RCC5 network, this yields all validity conditions for the
extracted quotient.

The result, combined with the soundness theorems
of~\cite{SplitForest}, gives: \emph{a concept $C_0$ in $\ALCIRCC{5}$
is satisfiable if and only if a valid finite rank-$d$ quotient for
$C_0$ exists}, where $d=\md(C_0)$.  Since the set of candidate
quotients is finite and validity is decidable, concept satisfiability
in $\ALCIRCC{5}$ is decidable.
\end{abstract}


% ═══════════════════════════════════════════════════════════════
\section{Introduction}
\label{sec:intro}
% ═══════════════════════════════════════════════════════════════

The description logic $\ALCIRCC{5}$ extends $\mathcal{ALCI}$ with
the five jointly exhaustive, pairwise disjoint RCC5 relations
$\DR$, $\PO$, $\EQ$, $\PP$, $\PPI$ as concrete roles, under
strong $\EQ$ semantics ($\EQ$ is identity).  Wessel~\cite{Wessel2001,Wessel2002,Wessel2003}
introduced the logic and left decidability open.

A recent split-forest proposal~\cite{SplitForest} attacks the problem
by decomposing models into cover trees (oriented by immediate $\PPI$
downward) and discharging the non-tree relations $\DR/\PO$ via finite
sibling-interface descriptors.  The proposal proves two main results:

\begin{enumerate}
\item \textbf{Finite-index lemma:} For each $k$, the set $\Sigma_k$
  of rank-$k$ states is finite and computable.
\item \textbf{Soundness:} Every valid finite rank-$d$ quotient unfolds
  into a genuine strong-$\EQ$ model of $C_0$.
\end{enumerate}

The soundness chain is fully detailed: valid quotient $\to$ unfolding
$\to$ disjunctive RCC5 network $\to$ arc-consistency (by LCA
induction) $\to$ K\"onig's lemma $\to$ canonical extendible
refinements $\to$ weak-$\EQ$ model $\to$ strong-$\EQ$ model.

\medskip
The \emph{completeness} direction---that every satisfiable concept
$C_0$ admits a valid finite rank-$d$ quotient---is stated as a proof
sketch in~\cite{SplitForest}, with the remark that ``the only
compressed part \ldots\ is the extraction of the support-closed tables
and relation-sensitive witness-menu filters from a model.''

The purpose of this note is to write out this extraction explicitly.
We show:

\begin{theorem}[Completeness extraction]\label{thm:main}
If\/ $\mathcal{I},d_\ast\models C_0$ in $\ALCIRCC{5}$ under
strong $\EQ$ semantics, then there exists a valid finite rank-$d$
quotient $Q$ for $C_0$ with $d=\md(C_0)$.
\end{theorem}

Combined with the soundness theorems of~\cite{SplitForest}, this
gives decidability of concept satisfiability in $\ALCIRCC{5}$.


% ═══════════════════════════════════════════════════════════════
\section{Preliminaries}
\label{sec:prelim}
% ═══════════════════════════════════════════════════════════════

We work under the definitions of~\cite{SplitForest} throughout.
We recall only the key notions needed for the extraction; see
\cite{SplitForest} for full details.

\subsection{Types and closure}

Fix a concept $C_0$ in negation normal form.  The \emph{Fischer-Ladner
closure} $\cl(C_0)$ is the standard closure under subformulas and
negation normal form duals.  A \emph{Hintikka type}
$\tau\in\Tp(C_0)\subseteq 2^{\cl(C_0)}$ is a maximally consistent
subset satisfying the usual propositional closure conditions.
For an element $d$ in a model~$\mathcal{I}$, its \emph{type}
$\tp(d)=\{D\in\cl(C_0):\mathcal{I},d\models D\}$ is a Hintikka type.

\subsection{Per-relation need families}

\begin{definition}[Need families, {\cite[Definition~1.4]{SplitForest}}]
\label{def:need}
For a Hintikka type $\tau\in\Tp(C_0)$ and a relation
$R\in\RCC$, define
\[
\Need_R(\tau) := \{\,D\in\cl(C_0) \mid \forall R.D\in\tau\,\}.
\]
A relation $S\in\RCC$ on an edge from type $\tau_1$ to type $\tau_2$
is \emph{relation-safe} if
$\Need_S(\tau_1)\subseteq\tau_2$ and
$\Need_{\inv(S)}(\tau_2)\subseteq\tau_1$.
\end{definition}

\begin{lemma}[Model relations are relation-safe]\label{lem:model-safe}
In any model $\mathcal{I}$ of $\ALCIRCC{5}$, if $\rho(d_1,d_2)=S$
for distinct elements $d_1,d_2$, then $S$ is relation-safe for the
pair $(\tp(d_1),\tp(d_2))$.
\end{lemma}

\begin{proof}
Let $D\in\Need_S(\tp(d_1))$, i.e., $\forall S.D\in\tp(d_1)$, i.e.,
$\mathcal{I},d_1\models\forall S.D$.  Since $\rho(d_1,d_2)=S$, the
semantics of $\forall S$ gives $\mathcal{I},d_2\models D$, hence
$D\in\tp(d_2)$.  Thus $\Need_S(\tp(d_1))\subseteq\tp(d_2)$.

The dual condition $\Need_{\inv(S)}(\tp(d_2))\subseteq\tp(d_1)$
follows by the same argument using $\rho(d_2,d_1)=\inv(S)$.
\end{proof}

This elementary observation is the \emph{key lemma} for the entire
extraction.  It ensures that when we quotient a model by rank-$d$
state equality, every relation that appears in the resulting
disjunctive domains is already relation-safe.

\subsection{Rank-$d$ states}

A \emph{rank-$0$ state} is a Hintikka type $\tau\in\Tp(C_0)$.
A \emph{rank-$(k{+}1)$ state} of a tree node $u$ is the tuple
\[
\st_{k+1}(u) = \bigl(\lambda(u),\;\Par_k(u),\;\Child_k(u),\;
  \Iface_k(u),\;\Wit_k(u)\bigr),
\]
recording the local type, parent's rank-$k$ state, child multiplicities
(in $\{0,1,2{+}\}$ per child rank-$k$ state), sibling-interface
descriptors between ordered child pairs, and a finite witness menu.
A \emph{rank-$d$ state} is an element of $\Sigma_d$ with $d=\md(C_0)$.
By the finite-index lemma~\cite[Theorem~1.25]{SplitForest}, each
$\Sigma_k$ is finite.

\subsection{Sibling-interface descriptors}

For an ordered sibling pair $(s,t)$ with $\rho(s,t)=\PO$, each
descendant~$x$ of~$s$ is classified by the three-way status partition
into $\Core(s,t)$, $\Out(s,t)$, or $\Front(s,t)$; see
\cite[Definition~1.6, Theorem~1.4]{SplitForest}.

The \emph{rank-$k$ sibling-interface descriptor} for $(s,t)$ consists of:
\begin{itemize}
\item Core multiplicities $K_k(s,t)$;
\item Open support sets $L_k(s,t)$, $R_k(s,t)$ of endpoint summaries
  (pairs of rank-$k$ state and status);
\item A relation table $M_k(s,t)$ mapping endpoint-summary pairs to
  disjunctive domains in
  $\mathbb{D}=\{\{\DR\},\{\PO\},\{\DR,\PO\}\}$.
\end{itemize}

A descriptor is \emph{locally coherent} if it satisfies
conditions~(C1)--(C5) of~\cite[Definition~1.8]{SplitForest}:
\begin{enumerate}[label=(C\arabic*)]
\item $\Out$ endpoints yield $\{\DR\}$ (DR rigidity).
\item $\Front$--$\Front$ entries are type-safe subsets of
  $\{\DR,\PO\}$.
\item Left downward support: child-step evolution preserves table
  entries via $\PP$-composition.
\item Right downward support: symmetric to~(C3).
\item Triple coherence: three-child consistency via RCC5 composition.
\end{enumerate}

\subsection{Valid finite quotient}

A \emph{valid finite rank-$d$ quotient}
$Q=(S,\sigma^\uparrow,\sigma_\ast,\ldots)$
consists of a finite set $S\subseteq\Sigma_d$ of rank-$d$ states, a
distinguished ambient-root state $\sigma^\uparrow$, a distinguished
evaluation state $\sigma_\ast$, and the inherited parent/child/interface/witness
data, satisfying~\cite[Definition~1.10]{SplitForest}:
\begin{enumerate}[label=(V\arabic*)]
\item Every existential in each state's type is matched by the witness
  menu.
\item Every universal is respected by the type-safe relation domains.
\item Every explicit witness-menu edge and every \emph{induced}
  pair-domain (from composition with tree edges) has non-empty
  $\Need_R$-filtered domain.
\item Repeated states are quotiented by rank-$d$ equality.
\end{enumerate}


% ═══════════════════════════════════════════════════════════════
\section{The extraction}
\label{sec:extraction}
% ═══════════════════════════════════════════════════════════════

Let $\mathcal{I}$ be a model of $\ALCIRCC{5}$ with
$\mathcal{I},d_\ast\models C_0$.  We construct a valid finite
rank-$d$ quotient in five steps.

\subsection{Step~1: Split-tree presentation}

\begin{construction}[Split tree]\label{con:split-tree}
Let $\Delta$ denote the domain of $\mathcal{I}$ and
$\rho:\Delta\times\Delta\to\{\DR,\PO,\EQ,\PP,\PPI\}$ the RCC5
labeling (with $\rho(d,d)=\EQ$ and strong $\EQ$ semantics).

\begin{enumerate}
\item \textbf{PP-DAG.}\; Define the proper-part DAG
  $D=(\Delta,\le_{\PP})$ where $d_1\le_{\PP} d_2$ iff
  $\rho(d_1,d_2)=\PP$.  Since $\PP$ is a strict partial order
  ($\comp(\PP,\PP)=\{\PP\}$, irreflexive by strong $\EQ$), this is
  a well-founded DAG.

\item \textbf{Ambient root.}\;
  If $d_\ast$ has no $\PP$-predecessor, set $d_\ast$ as the root of
  the DAG.  Otherwise, choose any $\le_{\PP}$-minimal element as the
  ambient root; if no unique root exists, add a fresh element $d_\top$
  with $\rho(d_\top,d)=\PP$ for all $d\in\Delta$ and
  $\tp(d_\top)=\tau_\top$ for a suitable type
  $\tau_\top\supseteq\{\top\}$.

\item \textbf{Tree unfolding.}\;
  Unfold the DAG into a rooted tree $T$ by the standard tree-unfolding
  of the immediate-$\PPI$ relation (Hasse diagram of $\le_{\PP}$).
  Each tree node $u\in T$ has a unique path from the root, and
  $\orig(u)\in\Delta$ maps it to the corresponding domain element.
  Join nodes (elements with multiple immediate $\PP$-predecessors)
  produce multiple tree copies; distinct copies of the same element
  are $\EQ$-mates.

\item \textbf{Type assignment.}\; For each $u\in T$, set
  $\lambda(u)=\tp(\orig(u))$.
\end{enumerate}
\end{construction}

\begin{remark}
The tree unfolding may be infinite (e.g., when the model has an
infinite $\PP$-chain).  This is not a problem: we will quotient by
rank-$d$ state equality, and only finitely many states arise.
\end{remark}

\subsection{Step~2: Rank-$d$ states from the model}

\begin{construction}[Rank-$d$ states]\label{con:states}
Assign rank-$k$ states to tree nodes by induction on~$k$.
\begin{itemize}
\item \textbf{Rank~$0$:}\; $\st_0(u)=\lambda(u)\in\Tp(C_0)$.
\item \textbf{Rank~$k{+}1$:}\; For each node $u$, define
  $\st_{k+1}(u)=(\lambda(u),\Par_k(u),\Child_k(u),\Iface_k(u),\Wit_k(u))$
  where:
  \begin{enumerate}[label=(\alph*)]
  \item $\Par_k(u)=\st_k(\mathsf{par}(u))$ if $u$ has a parent,
    $\bot$ if $u$ is the ambient root;
  \item $\Child_k(u):\Sigma_k\to\{0,1,2{+}\}$ records the multiplicity
    of each rank-$k$ state among $u$'s children;
  \item $\Iface_k(u)$ is the family of rank-$k$ sibling-interface
    descriptors between ordered child pairs, \emph{extracted from the
    model's realized relations} (Construction~\ref{con:descriptors});
  \item $\Wit_k(u)$ is the witness menu \emph{extracted from the
    model's existential witnesses} (Construction~\ref{con:witness}).
  \end{enumerate}
\end{itemize}
\end{construction}

\subsection{Step~3: Extracting sibling-interface descriptors}

\begin{construction}[Descriptor extraction]\label{con:descriptors}
Let $(s,t)$ be an ordered sibling pair (children of the same parent)
in the split tree~$T$, with $\rho(\orig(s),\orig(t))\in\{\DR,\PO\}$
(the only possibilities for distinct siblings after $\EQ$-splitting).

If $\rho(\orig(s),\orig(t))=\DR$, the descriptor is trivial: all
cross-subtree relations are $\DR$ (by $\comp(\PPI,\DR)=\{\DR\}$ and
$\comp(\DR,\PP)=\{\DR\}$).

If $\rho(\orig(s),\orig(t))=\PO$, construct the rank-$k$
descriptor as follows.

\begin{enumerate}
\item \textbf{Status assignment.}\; For each descendant $x$ of $s$
  (including $s$ itself), define $\mathsf{stat}(x)\in\{\Core,\Out,\Front\}$
  according to~\cite[Definition~1.6]{SplitForest}:
  \begin{itemize}
  \item $\Core$: $\orig(x)$ lies below both $\orig(s)$ and $\orig(t)$
    in the PP-DAG;
  \item $\Out$: $\orig(x)\notin\Core$ and $\rho(\orig(x),\orig(t))=\DR$;
  \item $\Front$: $\orig(x)\notin\Core$ and $\rho(\orig(x),\orig(t))=\PO$.
  \end{itemize}

\item \textbf{Core multiplicities.}\;
  $K_k(s,t)(\sigma) = \min(|\{x\in T_s : \st_k(x)=\sigma,\;
  \mathsf{stat}(x)=\Core\}|, 2)$ for each $\sigma\in\Sigma_k$.

\item \textbf{Open support sets.}\;
  \begin{align*}
  L_k(s,t) &= \{(\st_k(x),\mathsf{stat}(x)) :
    x\in T_s,\;\mathsf{stat}(x)\in\{\Out,\Front\}\}, \\
  R_k(s,t) &= \{(\st_k(y),\mathsf{stat}(y)) :
    y\in T_t,\;\mathsf{stat}(y)\in\{\Out,\Front\}\}.
  \end{align*}

\item \textbf{Relation table.}\;
  For $(p,\sigma)\in L_k(s,t)$ and $(q,\tau)\in R_k(s,t)$, define
  \[
  M_k(s,t)\bigl((p,\sigma),(q,\tau)\bigr) =
  \{\rho(\orig(x),\orig(y)) : \st_k(x)=p,\;\mathsf{stat}(x)=\sigma,\;
    \st_k(y)=q,\;\mathsf{stat}(y)=\tau\}.
  \]
  That is: the set of RCC5 relations that \emph{actually appear} in
  the model between elements whose rank-$k$ states and statuses match
  the given endpoint summaries.
\end{enumerate}
\end{construction}

\subsection{Step~4: Extracting witness menus}

\begin{construction}[Witness-menu extraction]\label{con:witness}
For each tree node $u$ and each existential
$\exists R.D\in\lambda(u)$, the model provides a witness:
an element $w\in\Delta$ with $\rho(\orig(u),w)=R$ and
$\mathcal{I},w\models D$.

\begin{itemize}
\item If $R\in\{\PP,\PPI\}$: the witness is an ancestor or descendant
  in the split tree.  These are handled by the tree structure
  (parent/child slots) and do not enter the witness menu.

\item If $R\in\{\DR,\PO\}$: the witness is a cross-edge target.
  Find a tree node $v$ with $\orig(v)=w$ (at least one exists since
  $\orig$ is surjective onto $\Delta$).  The witness-menu entry for
  this existential is:
  \[
  (\exists R.D) \;\mapsto\; (\st_k(v),\; R).
  \]
\end{itemize}

The witness menu $\Wit_k(u)$ is the collection of all such entries
for the $\DR/\PO$-existentials in $\lambda(u)$.
\end{construction}

\subsection{Step~5: Forming the quotient}

\begin{construction}[Quotient]\label{con:quotient}
Define the finite rank-$d$ quotient $Q$ as follows.
\begin{itemize}
\item $S = \{\st_d(u) : u\in T\} \subseteq \Sigma_d$.
  By the finite-index lemma, $|S|\le|\Sigma_d|<\infty$.
\item $\sigma^\uparrow = \st_d(\varrho)$ where $\varrho$ is the
  ambient root of $T$.
\item $\sigma_\ast = \st_d(u_\ast)$ where $u_\ast$ is the tree node
  corresponding to $d_\ast$ (the evaluation point).
\item Parent/child/interface/witness data are inherited from the
  state definitions (Construction~\ref{con:states}).
\end{itemize}
\end{construction}


% ═══════════════════════════════════════════════════════════════
\section{Verification of local coherence}
\label{sec:coherence}
% ═══════════════════════════════════════════════════════════════

We verify that the extracted sibling-interface descriptors satisfy
the local coherence conditions (C1)--(C5)
of~\cite[Definition~1.8]{SplitForest}.

\begin{lemma}[Coherence of extracted descriptors]\label{lem:coherence}
The descriptor extracted in Construction~\ref{con:descriptors}
satisfies conditions \emph{(C1)--(C5)}.
\end{lemma}

\begin{proof}
Let $(s,t)$ be an ordered sibling pair with
$\rho(\orig(s),\orig(t))=\PO$.

\medskip\noindent\textbf{(C1): Out rigidity.}\;
If $\mathsf{stat}(x)=\Out$, then $\rho(\orig(x),\orig(t))=\DR$.
For any descendant $y$ of $t$ with $\mathsf{stat}(y)\in\{\Out,\Front\}$:
$\rho(\orig(x),\orig(y))\in\comp(\DR,\PPI^*)$, and by
$\comp(\DR,\PPI)=\{\DR\}$ we get $\rho(\orig(x),\orig(y))=\DR$.
Hence the table entry is $\{\DR\}$.\;\checkmark

\medskip\noindent\textbf{(C2): Type-safety.}\;
Let $(p,\Front)\in L_k$ and $(q,\Front)\in R_k$, and let
$R\in M_k((p,\Front),(q,\Front))$.  Then $R$ is realized by some pair
$(x,y)$ with $\st_k(x)=p$, $\st_k(y)=q$, $\rho(\orig(x),\orig(y))=R$.
By Lemma~\ref{lem:model-safe}, $R$ is relation-safe for
$(\lambda(x),\lambda(y))=(\lambda(p),\lambda(q))$.
Hence $M_k((p,\Front),(q,\Front))\subseteq\Safe(p,q)\cap\{\DR,\PO\}$.
Moreover, $M_k\ne\varnothing$ since at least one pair $(x,y)$
realizes it.
\;\checkmark

\medskip\noindent\textbf{(C3): Left downward support.}\;
Let $(p,\sigma)\in L_k$ evolve to $(p',\sigma')$ via a child step
(i.e., $x$ has a child $x'$ with $\st_k(x')=p'$ and
$\mathsf{stat}(x')=\sigma'$).  For any $(q,\tau)\in R_k$ and any
$R\in M_k((p,\sigma),(q,\tau))$, there exist witnesses $x,y$ with
$\rho(\orig(x),\orig(y))=R$.

Since $x'$ is a child of $x$ in the split tree,
$\rho(\orig(x),\orig(x'))=\PPI$ (immediate parent--child).  By the
RCC5 composition table:
\[
\rho(\orig(x'),\orig(y)) \;\in\;
  \comp(\PP,\rho(\orig(x),\orig(y))) \;=\; \comp(\PP,R).
\]
Let $R'=\rho(\orig(x'),\orig(y))$.  Then $R'\in\comp(\PP,R)$ and
$R'\in M_k((p',\sigma'),(q,\tau))$ (since $x'$ witnesses the left
endpoint summary and $y$ witnesses the right).
\;\checkmark

\medskip\noindent\textbf{(C4): Right downward support.}\;
Symmetric to~(C3).\;\checkmark

\medskip\noindent\textbf{(C5): Triple coherence.}\;
Let $(s_1,s_2,s_3)$ be three siblings and let $x_i\in T_{s_i}$ for
$i=1,2,3$ with endpoint summaries $\chi_i$.  The relation table
entries $M((\chi_1,\chi_2))$ and $M((\chi_1,\chi_3))$ are extracted
from the model.  For any $R_{12}\in M((\chi_1,\chi_2))$ and
$R_{13}\in M((\chi_1,\chi_3))$, there exist witnesses
$x_1,x_2,x_3$ realizing these relations.  By the model's
composition-consistency:
\[
\rho(\orig(x_2),\orig(x_3)) \;\in\;
  \comp(\inv(R_{12}),R_{13}).
\]
This relation appears in $M((\chi_2,\chi_3))$, so the triple
coherence condition is satisfied by the model's own atomic
consistency.
\;\checkmark
\end{proof}


% ═══════════════════════════════════════════════════════════════
\section{Relation-safety of extracted witness menus}
\label{sec:relation-safety}
% ═══════════════════════════════════════════════════════════════

\begin{lemma}[Witness-menu relation-safety]\label{lem:witness-safe}
Every entry in the witness menu extracted in
Construction~\ref{con:witness} is relation-safe.
\end{lemma}

\begin{proof}
Let $\exists R.D\in\lambda(u)$ with $R\in\{\DR,\PO\}$, and let
$w$ be the model witness with $\rho(\orig(u),w)=R$ and
$\mathcal{I},w\models D$.  Let $v$ be a tree node with
$\orig(v)=w$.  The witness-menu entry is $(\st_k(v), R)$.

By Lemma~\ref{lem:model-safe}, $R$ is relation-safe for
$(\tp(\orig(u)),\tp(w))=(\lambda(u),\lambda(v))$.

Since $\lambda(u)$ and $\lambda(v)$ depend only on
$\st_k(u)$ and $\st_k(v)$ (the local type is the first component
of every rank-$k$ state), the relation-safety condition
\[
\Need_R(\lambda(u))\subseteq\lambda(v)
\quad\text{and}\quad
\Need_{\inv(R)}(\lambda(v))\subseteq\lambda(u)
\]
holds for the state pair $(\st_k(u),\st_k(v))$.
\end{proof}


% ═══════════════════════════════════════════════════════════════
\section{Non-emptiness of filtered domains}
\label{sec:nonempty}
% ═══════════════════════════════════════════════════════════════

The most delicate part of the validity verification is showing that
\emph{induced pair-domains}---those arising from composing
witness-menu edges with tree edges---have non-empty
$\Need_R$-filtered domains.

\begin{lemma}[Filtered domains are non-empty]\label{lem:filtered}
Let $Q$ be the quotient extracted in Construction~\ref{con:quotient}.
For any two states $\sigma_1,\sigma_2\in S$ such that some pair of
elements in the model realizes a relation between elements with
states $\sigma_1$ and $\sigma_2$, the $\Need_R$-filtered domain
\[
\Dom'(\sigma_1,\sigma_2) = \{R\in\Dom(\sigma_1,\sigma_2) :
  \Need_R(\sigma_1)\subseteq\lambda(\sigma_2)\text{ and }
  \Need_{\inv(R)}(\sigma_2)\subseteq\lambda(\sigma_1)\}
\]
is non-empty, where
$\Dom(\sigma_1,\sigma_2)=\{\rho(\orig(x),\orig(y)):
\st_d(x)=\sigma_1,\;\st_d(y)=\sigma_2\}$.
\end{lemma}

\begin{proof}
By assumption, there exist tree nodes $x,y$ with
$\st_d(x)=\sigma_1$, $\st_d(y)=\sigma_2$, and
$R=\rho(\orig(x),\orig(y))\in\RCC$.
By Lemma~\ref{lem:model-safe}, $R$ is relation-safe for
$(\tp(\orig(x)),\tp(\orig(y)))=(\lambda(\sigma_1),\lambda(\sigma_2))$.
Hence $R\in\Dom'(\sigma_1,\sigma_2)$.
\end{proof}

\begin{remark}[Filtered = unfiltered for model-extracted quotients]
\label{rem:filtered-eq-unfiltered}
Lemma~\ref{lem:filtered} shows that \emph{every} relation in the
unfiltered domain $\Dom(\sigma_1,\sigma_2)$ is already
relation-safe, since each such relation is witnessed by a pair
from the model.  Therefore
$\Dom'(\sigma_1,\sigma_2)=\Dom(\sigma_1,\sigma_2)$ for
model-extracted quotients.  The $\Need_R$-filtering step, which is
crucial for \emph{checking} candidate quotients, is vacuously
satisfied when \emph{extracting} a quotient from a genuine model.
\end{remark}

The above handles explicit witness-menu edges.  We now address
\emph{induced} pair-domains from composition.

\begin{lemma}[Composition-induced domains]\label{lem:composition}
Let $\sigma_1,\sigma_2,\sigma_3\in S$ be states such that elements
$x_1,x_2,x_3$ in the model realize
$\rho(\orig(x_1),\orig(x_2))=R_{12}$ with $\st_d(x_i)=\sigma_i$.
Then the induced domain on $(\sigma_1,\sigma_3)$ via composition
through $\sigma_2$ contains at least $\rho(\orig(x_1),\orig(x_3))$,
which is relation-safe.
\end{lemma}

\begin{proof}
By composition-consistency of the model:
$\rho(\orig(x_1),\orig(x_3))\in\comp(R_{12},\rho(\orig(x_2),\orig(x_3)))$.
Hence
$\rho(\orig(x_1),\orig(x_3))\in\Dom(\sigma_1,\sigma_3)$.
By Lemma~\ref{lem:model-safe}, this relation is relation-safe for
$(\lambda(\sigma_1),\lambda(\sigma_3))$.
Hence it survives $\Need_R$-filtering.
\end{proof}


% ═══════════════════════════════════════════════════════════════
\section{Validity of the extracted quotient}
\label{sec:validity}
% ═══════════════════════════════════════════════════════════════

\begin{theorem}[Validity]\label{thm:validity}
The quotient $Q$ extracted in Construction~\ref{con:quotient} is a
valid finite rank-$d$ quotient for $C_0$ in the sense
of~\cite[Definition~1.10]{SplitForest}.
\end{theorem}

\begin{proof}
We verify conditions (V1)--(V4).

\medskip\noindent\textbf{(V1): Existential satisfaction.}\;
For each $\sigma\in S$ and each $\exists R.D\in\lambda(\sigma)$:
\begin{itemize}
\item If $R\in\{\PP,\PPI\}$: the existential is an eventuality over
  the ancestor/descendant order of the cover tree.  In the model,
  a witness exists as an ancestor or descendant of any element with
  state $\sigma$.  By the construction of parent/child slots, this
  witness's state is reachable via parent or child navigation in~$Q$.
\item If $R\in\{\DR,\PO\}$: the witness is recorded in
  $\Wit_d(\sigma)$ by Construction~\ref{con:witness}, with target
  state in $S$ and relation $R$.  By Lemma~\ref{lem:witness-safe},
  the entry is relation-safe.
\end{itemize}

\medskip\noindent\textbf{(V2): Universal satisfaction.}\;
For each $\sigma\in S$ and each $\forall R.D\in\lambda(\sigma)$:
$D\in\Need_R(\lambda(\sigma))$.  For any $\sigma'\in S$ with
$R\in\Dom(\sigma,\sigma')$, there exist witnesses $x,y$ with
$\rho(\orig(x),\orig(y))=R$, $\st_d(x)=\sigma$, $\st_d(y)=\sigma'$.
Since $\mathcal{I},\orig(x)\models\forall R.D$ and
$\rho(\orig(x),\orig(y))=R$, we get $D\in\tp(\orig(y))=\lambda(\sigma')$.
Hence universals are respected.\;\checkmark

\medskip\noindent\textbf{(V3): Non-empty filtered domains.}\;
By Lemma~\ref{lem:filtered}, every pair-domain $\Dom(\sigma_1,\sigma_2)$
that arises in the quotient has non-empty $\Need_R$-filtered domain
(in fact, filtered $=$ unfiltered by
Remark~\ref{rem:filtered-eq-unfiltered}).  By
Lemma~\ref{lem:composition}, composition-induced domains also have
non-empty filtered domains.\;\checkmark

\medskip\noindent\textbf{(V4): Coherent descriptors.}\;
By Lemma~\ref{lem:coherence}, the extracted sibling-interface
descriptors satisfy (C1)--(C5).\;\checkmark

\medskip\noindent\textbf{Evaluation state contains $C_0$.}\;
The tree node $u_\ast$ corresponding to $d_\ast$ has
$\lambda(u_\ast)=\tp(d_\ast)\ni C_0$.  Hence
$C_0\in\lambda(\sigma_\ast)$.\;\checkmark

\medskip\noindent\textbf{Finiteness.}\;
$|S|\le|\Sigma_d|<\infty$ by the finite-index lemma.\;\checkmark
\end{proof}


% ═══════════════════════════════════════════════════════════════
\section{Main theorem}
\label{sec:main}
% ═══════════════════════════════════════════════════════════════

\begin{proof}[Proof of Theorem~\ref{thm:main}]
Combine Construction~\ref{con:quotient} with
Theorem~\ref{thm:validity}: the extracted quotient is valid, finite,
and has $C_0$ in its evaluation state's type.
\end{proof}

\begin{corollary}[Decidability of $\ALCIRCC{5}$]\label{cor:decidability}
Concept satisfiability in $\ALCIRCC{5}$ under strong $\EQ$ semantics
is decidable.
\end{corollary}

\begin{proof}
By the finite-index lemma~\cite[Theorem~1.25]{SplitForest}, the
number of candidate quotients (up to isomorphism) is finite.
Validity of a quotient is decidable (it involves checking finite
conditions on finite data).

\begin{itemize}
\item If $C_0$ is satisfiable, Theorem~\ref{thm:main} produces a
  valid quotient.  The enumeration will find one.
\item If a valid quotient exists, the soundness
  theorems~\cite[Theorems~1.18, 1.19]{SplitForest} produce a genuine
  model.  Hence $C_0$ is satisfiable.
\end{itemize}

Therefore: $C_0$ is satisfiable if and only if a valid finite rank-$d$
quotient for $C_0$ exists, and the search terminates.
\end{proof}


% ═══════════════════════════════════════════════════════════════
\section{Discussion}
\label{sec:discussion}
% ═══════════════════════════════════════════════════════════════

\subsection{Why the extraction is straightforward}

The key insight (Lemma~\ref{lem:model-safe}) is that in any valid
$\ALCIRCC{5}$ model, every realized relation between two elements is
already relation-safe for their types.  This follows directly from
the semantics of universal quantification: $\forall R.D$ means that
\emph{every} $R$-related element satisfies $D$.

When we quotient the model by rank-$d$ state equality, the
disjunctive domains $\Dom(\sigma_1,\sigma_2)$ consist of relations
witnessed by actual model pairs.  Each such relation is
relation-safe.  Therefore the $\Need_R$-filtering step---which is
the most complex part of the \emph{validity checking}
procedure---becomes vacuous for extracted quotients.

This observation was implicit in~\cite{SplitForest} (``the $\Need_R$
families are read off directly from local types'') but deserves to
be stated explicitly, since it is the reason the extraction works.

\subsection{The role of the patchwork property}

The patchwork property of RCC5 (path-consistent $\Rightarrow$
globally consistent) is used in the \emph{soundness}
direction~\cite{SplitForest}: it guarantees that the disjunctive
network produced by unfolding a valid quotient can be refined to an
atomic, globally consistent assignment.

In the \emph{completeness} direction (this note), the patchwork
property is \emph{not} directly needed.  The model already provides
a globally consistent atomic assignment; we are simply reading off
the quotient from it.  The composition-consistency of the model's
RCC5 network ensures coherence conditions (C3)--(C5) and the
non-emptiness of composition-induced domains
(Lemma~\ref{lem:composition}).

\subsection{Complexity}

The quotient search enumerates all subsets $S\subseteq\Sigma_d$
with associated data and checks validity.  Since $|\Sigma_d|$ is
bounded by a tower of exponentials in $|\cl(C_0)|$, the resulting
procedure is decidable but may not be elementary.

An $\textsc{ExpTime}$ upper bound would require showing that
$|\Sigma_d|$ is at most singly exponential in $|C_0|$.  This is an
\emph{independent open problem}: even the PO-coherent
fragment~\cite{TwoTier} achieves only a doubly exponential quotient
bound.

\subsection{Honest assessment}

We believe the argument is complete, but we flag the following
points for scrutiny:

\begin{enumerate}
\item \textbf{Ambient root choice.}\; The construction may require
  adding a fresh top element $d_\top$ if the model has no unique
  $\le_{\PP}$-minimal element.  The type $\tau_\top$ assigned to
  $d_\top$ must be a valid Hintikka type; this requires $\top\in\cl(C_0)$
  (which holds by standard closure conventions) and that $\tau_\top$
  is compatible with having every domain element as a $\PPI$-successor.
  In practice this is unproblematic, but a fully rigorous treatment
  would spell out the construction of $\tau_\top$.

\item \textbf{PP/PPI eventualities.}\; The verification of~(V1) for
  $\PP$/$\PPI$-existentials asserts that witnesses are reachable via
  parent/child navigation.  This is correct because the split tree
  faithfully represents the $\PP/\PPI$ partial order, but the
  precise argument involves the interaction between tree unfolding
  and rank-$d$ state quotient.  A witness $w$ with
  $\rho(\orig(u),w)=\PP$ is an ancestor of $u$ in the split tree;
  its state $\st_d(w)$ appears in the parent chain of $\sigma$, which
  is recorded in the quotient's parent-slot data.

\item \textbf{Support closure.}\; The witness-menu target state must
  lie in~$S$.  This is immediate from the construction: the target is
  $\st_d(v)$ for a tree node~$v$, which is in~$S$ by definition.
  But if the model is infinite and the quotient is finite, distinct
  model elements may map to the same state.  This is fine: the
  quotient identifies them, and the witness menu points to the
  equivalence class, not to a specific element.

\item \textbf{Induced pair-domains for deeper compositions.}\;
  Lemma~\ref{lem:composition} handles one level of composition.
  For deeper chains ($\sigma_1 \to \sigma_2 \to \sigma_3 \to \sigma_4$),
  the argument iterates: at each step, the model provides actual
  elements realizing the composed relation, and
  Lemma~\ref{lem:model-safe} applies.  The model's
  composition-consistency ensures that all multi-step compositions
  are realized.
\end{enumerate}

None of these points appear to harbor a genuine gap.  The extraction
rests on a single clean principle: \emph{model relations are
relation-safe}, and this is a direct consequence of the semantics of
universal quantification.


% ═══════════════════════════════════════════════════════════════
\section{Conclusion}
\label{sec:conclusion}
% ═══════════════════════════════════════════════════════════════

We have written out the completeness extraction that was compressed
in~\cite{SplitForest}.  The argument decomposes into:
\begin{enumerate}
\item Split-tree presentation (standard).
\item Rank-$d$ state assignment (standard induction).
\item Descriptor extraction from realized model relations
  (coherence by model consistency).
\item Witness-menu extraction from model witnesses
  (relation-safety by Lemma~\ref{lem:model-safe}).
\item Quotient formation (finiteness by the finite-index lemma).
\item Validity verification (all conditions follow from model
  validity and the key lemma).
\end{enumerate}

Combined with the soundness chain of~\cite{SplitForest}
(Theorems~1.18 and~1.19), this yields:

\begin{center}
\emph{Concept satisfiability in $\ALCIRCC{5}$ under strong $\EQ$
semantics is decidable.}
\end{center}

The exact complexity is open.  The quotient search is decidable but
potentially non-elementary; an $\textsc{ExpTime}$ bound (matching
the known $\mathcal{ALCI}$ lower bound) would require a tighter
analysis of the rank-$d$ state space.

\medskip
\noindent\textbf{Computational evidence.}\;
An independent implementation~\cite{ClaudeImpl} cross-validates
the cover-tree tableau (which is the operational counterpart of the
split-forest quotient search) against a quasimodel-based reasoner
on 911 test concepts with zero mismatches.  A decomposition test
confirms that 775/775 satisfiable concepts (100\%) have finite
cover-tree models across 12.2\,M enumerated models, with zero
genuine counterexamples.  This provides strong empirical support for
the theoretical result.


\begin{thebibliography}{9}

\bibitem{Wessel2001}
M.~Wessel.
\newblock Obstacles on the way to qualitative spatial reasoning with
  description logics.
\newblock In \emph{Proc.\ DL Workshop}, 2001.

\bibitem{Wessel2002}
M.~Wessel.
\newblock On spatial reasoning with description logics---position paper.
\newblock In \emph{Proc.\ DL Workshop}, 2002.

\bibitem{Wessel2003}
M.~Wessel.
\newblock Qualitative spatial reasoning with the
  {$\mathcal{ALCI}_\mathit{RCC}$} family---{F}irst results and
  unanswered questions.
\newblock Technical Report FBI-HH-M-324/03, University of Hamburg, 2002/2003.

\bibitem{SplitForest}
GPT-5.4 (OpenAI).
\newblock Finite sibling-interface descriptors and a proposed completion
  for $\mathcal{ALCI}_{RCC5}$.
\newblock Proof-style note, April 2026.
\newblock Prompted by Michael Wessel.

\bibitem{TwoTier}
Claude (Anthropic).
\newblock Two-tier quotient decidability for the PO-coherent fragment
  of $\mathcal{ALCI}_{RCC5}$.
\newblock Proof-style note, March 2026.
\newblock Prompted by Michael Wessel.

\bibitem{ClaudeImpl}
Claude (Anthropic).
\newblock Cover-tree tableau implementation for $\mathcal{ALCI}_{RCC5}$.
\newblock Implementation paper, April 2026.
\newblock Prompted by Michael Wessel.

\end{thebibliography}

\end{document}
